% ============================================================
% f2014_liaoning_arts_q22_circle.tex
% 2014 辽宁卷（文科）第 22 题：圆的切线与割线
% 来源 PDF：普通高考/2014/2014辽宁文.pdf
% 原页：第 2 页（A3），印面 898，第 22 题题干下方
% EP 交圆于 E、C，PD 切圆于 D，G 在 CE 上且 PG=PD，
% DG 延长交圆于 A，AB 垂直 EP 于 F。
% 取单位圆：E(-0.866,-0.5)、C(0.866,-0.5)、A(0,-1)、B(0,1)、
% D(0.574,0.819)、P(2.456,-0.5)，此时 PD 与圆相切、PG=PD。
% ============================================================
\begin{tikzpicture}[>=Stealth, line width=0.7pt, font=\small,
  x=1.5cm, y=1.5cm, every node/.style={inner sep=1.5pt}]
  \coordinate (E) at (-0.866,-0.5);
  \coordinate (C) at (0.866,-0.5);
  \coordinate (A) at (0,-1);
  \coordinate (B) at (0,1);
  \coordinate (F) at (0,-0.5);
  \coordinate (G) at (0.158,-0.5);
  \coordinate (D) at (0.574,0.819);
  \coordinate (P) at (2.456,-0.5);

  % 圆
  \draw (0,0) circle (1);

  % 割线 EP、直径 AB、弦 BD、切线 PD、弦 ED、直线 DGA、弦 AC
  \draw (E) -- (P);
  \draw (A) -- (B);
  \draw (B) -- (D);
  \draw (D) -- (P);
  \draw (E) -- (D);
  \draw (D) -- (A);
  \draw (A) -- (C);

  % 点与标签
  \fill (E) circle (1.1pt) node[left] {$E$};
  \fill (C) circle (1.1pt) node[below] {$C$};
  \fill (A) circle (1.1pt) node[below] {$A$};
  \fill (B) circle (1.1pt) node[above] {$B$};
  \fill (F) circle (1.0pt) node[above left] {$F$};
  \fill (G) circle (1.0pt);
  \node at (0.24,-0.62) {$G$};
  \fill (D) circle (1.1pt) node[above right] {$D$};
  \fill (P) circle (1.1pt) node[right] {$P$};
\end{tikzpicture}
